\documentclass[11pt]{article}
\usepackage[T1]{fontenc}
\usepackage{lmodern}
\usepackage{amsmath,amssymb}
\usepackage[margin=0.9in]{geometry}
\usepackage[hidelinks]{hyperref}
\setlength{\parskip}{0.45em}
\setlength{\parindent}{0pt}
\newcommand{\loc}[2]{\textit{#1, lines #2.}}
\newcommand{\mt}[1]{\loc{\texttt{main.tex}}{#1}}
\newcommand{\ed}[1]{\loc{\texttt{explicit\_dynamic.tex}}{#1}}

\begin{document}

{\LARGE\bfseries Proof review}\\[0.3em]
{\large Consensus Clustering below Two}

\section*{Scope and assessment}

I reviewed the online, offline, implicit fully dynamic, and explicit
insertion-only results in the current \texttt{main.tex}, including the
explicit algorithm and its macros. Line numbers refer to the draft as
reviewed.

The online history bound and per-vertex recourse analysis check out with
the cited Online Agreement guarantees. I verified the offline charging
constants, Bernstein coefficients, and boundary factorization in exact
arithmetic. I found no substantive gap in the offline approximation or
the implicit fully dynamic maintenance argument. The explicit consensus
reduction also checks out conditional on its correlation-clustering solver.
The draft leaves the approximation and expected update time of that solver
unproved for the modified Dynamic Agreement algorithm.

\section{Gaps affecting the explicit dynamic theorem}

\subsection{Approximation and sparse anchor density need a proof for the modified algorithm}

\ed{48--59, 374--398, 699--707, 990--996}
The lemma for \textsc{ExplicitDynamicAgreement} promises a simultaneous
constant-factor approximation with failure probability at most \(N^{-2}\).
The next sections prove that \textsc{LabelSync} recovers connected
components conditional on a success event. They give neither a definition
and probability bound for that event nor an approximation proof for the
sparse graph the modified algorithm produces.

The published Dynamic Agreement analysis uses its own anchor updates.
The draft adds the \(L_t\) pruning rule, deletes edges after anchor
removals, and resets neighborhood snapshots. These operations change the
state used by subsequent \textsc{Clean}, \textsc{Anchor}, and
\textsc{Connect} calls. In particular, an argument that a sampled anchor
survives until the target prefix must account for the extra \(L_t\) test.
The draft also uses \(h_u\) without specifying which test sets it or when
it is refreshed. The precision boost for the tests appears only in a
comment at line 50. The source's Theorem~6 states a \(1-5/n\) bound and its
proof discusses a union over times and clusters. The draft needs to
justify the \(N^{-2}\) failure bound for its own algorithm. The consensus
reduction itself needs only probability \(1/2\) for simultaneous success.

The sparse-degree lemma is a separate dependency. Theorem~39 bounds
degrees in the input graph. \textsc{LabelSync} needs
\(\deg_{\widetilde G_t}(a)\ge 3|C|/4-1\) for every anchor \(a\)
in a nontrivial sparse component \(C\). Input density alone cannot give
this inequality. For example, the input could be a clique and its sparse
subgraph a spanning path. A possible route for the original algorithm
is to combine the sparse-neighborhood retention bounds in Lemma~33
with the stronger \((1-574434\varepsilon)|C|\) bound in the proof of
Theorem~39 at the anchor's last interesting event. If the draft keeps
the modified algorithm, it must prove the corresponding bounds for its
pruning and snapshot rules. Both approximation and sparse density need
to hold on the event used by \textsc{LabelSync}.
\href{https://openreview.net/attachment?id=3YG55Lbcnr&name=pdf}{Cohen-Addad et al.,
ICML 2024, Algorithms 3--6, Lemma 33, and Theorems 17, 23, 39, and 6.}

\subsection{The additional cleanup work is missing from the time bound}

\ed{374--398, 932--984}
The pseudocode closes the arrival handler before computing \(L_t\) and
calling \textsc{LabelSync}. Those operations need to run at every arrival.
Line 391 also says to delete \emph{vertices} with no endpoint in the
anchor set. The stated edge invariant suggests deleting \emph{edges}
with neither endpoint in the anchor set and retaining all active vertices.

Charging every edge deletion to a preceding insertion bounds the number
of deletions and leaves the cost of finding those edges unaccounted for. A full
scan of the sparse graph each update would take \(\Omega(n^2)\) total
time on successful executions on a disjoint-matching stream. Once a
positive pair in that stream has arrived, its edge must be retained in
the sparse graph to obtain zero disagreements. A local cleanup needs to
specify which edges it inspects when anchors change and charge those
inspections, including inspections that leave edges in place. The
analysis must also account for the degree checks and snapshot
maintenance. The current proof invokes (DA4), (B2), and (B3), which
are undefined in the draft. The modified trajectory needs its own
test-count bound.
\href{https://openreview.net/attachment?id=3YG55Lbcnr&name=pdf}{Cohen-Addad et al.,
ICML 2024, Appendix C.}

\section{Local proof and statement corrections}

\subsection{Distinguish the graph after deletions from the final graph}

\ed{480--518, 709--739, 751--778}
The robustness lemma establishes that vertices stranded from the old
anchors are singletons in \(\widetilde G_{t-1}[C]\setminus D_C\).
The correctness proof later asserts they are singletons in
\(\widetilde G_t\). Added edges can reconnect them to each other or to
other classes. For example, remove one leaf edge from a five-vertex old
star and attach that leaf to a newly arrived vertex. The leaf is a
singleton after deletions and belongs to a two-vertex component in the
final graph. Similarly, an old class with \(d_C=0\) remains internally
connected but can merge with another class. The described implementation
can handle both cases. In the few-deletions case, prove that every vertex
of \(C\setminus W_C\) remains connected to the anchor representative
\(r\) via surviving old edges. Use final connectivity queries to place
the vertices of \(W_C\), as the pseudocode does.

\subsection{Count the work of \textsf{LabelSync} on all executions}

\ed{584--602, 658--679, 784--785, 885--928, 956--975}
The statement of the per-update time lemma omits the \(+1\) already
present in its proof. An isolated arrival can have an empty journal
and zero changes to old labels, yet still requires processing and
reporting. Include the new assignment separately and use, for example,
\(O((|\mathcal J(t)|+1)\log(n+2)+c_t)\) when all vertices enumerated
in the merge step actually change labels.

For the unconditional expected-time proof, the paper should charge
enumeration to the vertices actually visited. The claim at line 920
that a piece is traversed only if its vertices receive a different label
uses the component invariant. On a failed density event, the retained
piece can contain vertices from different final components. If its key
is computed from an arbitrary vertex, a piece extracted from that same
old class can be merged back into the retained piece and enumerated
with its old label. Charge enumeration directly to the sizes of the
nonwinning pieces and the fresh-label winning pieces, plus one for the
arrival. The potential computation already bounds these sizes
irrespective of component correctness. When creating the report, compare
the final label with the stored old label so that an unchanged old
vertex is not reported.

\subsection{Specify the neighborhood convention for the Agreement tests}

\ed{34--46, 374--396}
The bucket implementation gives open neighborhoods. The draft switches
to closed neighborhoods for the pruning test and does not specify the
convention for agreement and heaviness tests. Using open neighborhoods
directly in the agreement test fails even on a graph with two vertices
and one positive edge. The endpoints have disjoint open neighborhoods
and fail the agreement threshold for any \(\varepsilon_{\mathrm{DA}}<1\),
although the optimum cost is zero. Spell out the closed-neighborhood
implementation where it is needed. A closed-neighborhood sample can be
generated by sampling the query vertex with probability \(1/(d+1)\)
and otherwise sampling an open neighbor. This also specifies the case
\(d=0\).

\subsection{State the prefix expectation used for the deterministic cap}

\ed{17--31, 286--327}
The reduction and the optional deterministic-cap proof use
\(\mathbb E W(u)\le uT_{\mathcal D}(n)\) for every \(u\le n\), with
the solver initialized using the same horizon \(n\). The solver lemma at
lines 48--52 promises this prefix bound, but the reduction's hypothesis
only says \emph{expected amortized update time}. State the bound explicitly
in the hypothesis so that both the prefix work proof and the dyadic
Markov argument have the premise they need.

\subsection{Align the theorem statements with the proved parameters}

\mt{201--221, 225--245, 445--461, 589--604, 1338--1356, 2368--2383}
The online theorem should quantify \(\varepsilon\) with the range
\(0<\varepsilon<2/(3\eta+5)\) used in the proof. Its \(\eta\) is the
competitive factor of the chosen deterministic Online Agreement
implementation with the required per-vertex recourse. The cited
theorems establish this guarantee without identifying \(\eta\) as the
best known factor among online correlation-clustering algorithms.

Use the positive logarithms from the proofs in the offline and
implicit time bounds, for example \(\log(N+2)\) instead of \(\log N\).
For the implicit query guarantee, define \(n=|V_t|\), use
\(O(\log(n+2))\) for lookup and \(O(n+1)\) for listing, and specify
whether listing returns distinct cluster identifiers or a vertex-label
pair for every active vertex. The proof currently describes the latter.
\href{https://proceedings.mlr.press/v162/cohen-addad22a.html}{Cohen-Addad et al.,
ICML 2022, Theorems E.2 and F.1.}

\end{document}
